% !PRINTED 427
\weberpage{427}
%<<CONT
die Congruenz identisch ist, d.\,h. dass alle Co\"efficienten von $f(x)$
durch $n$ theilbar sein m\"ussen.

Dieser Satz ist, wie man sieht, ganz analog dem algebraischen Satz, dass eine Gleichung nicht mehr Wurzeln haben
kann, als ihr Grad angiebt. Es l\"asst sich aber nicht der andere
Satz \"ubertragen, dass jede Gleichung auch wirklich so viele
Wurzeln hat. Eine Congruenz $m^{\text{ten}}$ Grades kann weniger, selbst
gar keine Wurzeln haben. Um so bemerkenswerther ist eine
besondere Congruenz, bei der die Zahl der Wurzeln immer dem
Grade gleichkommt, auf Grund eines Lehrsatzes, der der \gesperrt{Fermat'sche Lehrsatz} genannt wird, und den wir hier so formuliren.

\gesperrt{2. Die Congruenzen}
\weq{5}{x^n - x \equiv 0, \quad x^{n-1} - 1 \equiv 0 \quad (\mathrm{mod}\ n)}
\gesperrt{haben, wenn $n$ eine Primzahl ist, so viele Wurzeln
als ihr Grad betr\"agt, n\"amlich $n$ und $n-1$.}

Beide Behauptungen sind nicht wesentlich verschieden, denn
die erste der Congruenzen (5) hat alle Wurzeln der zweiten und
ausserdem die Wurzel 0, die der zweiten nicht gen\"ugt. Ebenso
hat die zweite alle Wurzeln der ersten, mit Ausnahme der
Wurzel 0.

Da es nun f\"ur den Modul $n$ \"uberhaupt nur $n$ verschiedene
Zahlen giebt, so ist also zu beweisen, dass f\"ur jede ganze Zahl
$\alpha$ die Congruenz besteht
\weq{6}{\alpha^n \equiv \alpha \quad (\mathrm{mod}\ n).}
Diese Congruenz ist richtig f\"ur $\alpha = 0$ und $\alpha = 1$. Wir
beweisen sie also wieder durch vollst\"andige Induction, indem wir
aus der als richtig vorausgesetzten Congruenz (6) die Richtigkeit von
\weq{7}{(\alpha + 1)^n \equiv \alpha + 1 \quad (\mathrm{mod}\ n)}
ableiten. Dies ist aber aus dem binomischen Satz zu schliessen,
wenn man beachtet, dass alle Binomialco\"efficienten, mit Ausnahme des ersten und des letzten, die gleich 1 sind, n\"amlich
\[
n, \quad \frac{n(n-1)}{1\,.\,2}, \quad \frac{n(n-1)(n-2)}{1\,.\,2\,.\,3}, \dots
\]
durch $n$ theilbare ganze Zahlen sind, weil der Z\"ahler, nicht aber
der Nenner durch $n$ theilbar ist. Demnach ist
\[
(\alpha + 1)^n \equiv \alpha^n + 1 \quad (\mathrm{mod}\ n),
\]
also folgt die Formel (7) aus der Formel (6).
